\chapter[Sermon 61. Of the Mark and Number of the Name of the Beast.]{Of the Mark and Number of the Name of the Beast.}
\label{sermon:061}
\markright{Sermon 61. Of the Mark and Number of the Name of the Beast.}

And he made all people, both the humble and the powerful, the wealthy and the impoverished, the free and the enslaved, to receive a mark on their right hands or on their foreheads. And that no one might buy or sell, except he who has the mark, or the name of the beast, or the number of his name. Here is wisdom. Let him who has understanding calculate the number of the beast, for it is a man’s number, and his number is six hundred threescore and six.

\index{Antichrist}He adds the rest, by which Antichrist may be recognized and avoided. Indeed, he may chiefly be known by these things that follow.

\index{Rome}And he speaks of the subjects of Antichrist and of this new king and bishop. He will procure to himself, says he, an infinite multitude of all kinds of men, of all states and degrees. For his kingdom shall be ample and large. Therefore the Lord resists here certain kinds and states of men: and under the same he understands whatsoever is of the same state in the whole world. The Roman Antichrist brought under his subjection small and great, rich and poor: free, namely nobles, and bond. For we see that emperors, kings, dukes, marquesses, earls and baronesses, realms, countries, cities, patriarchs, archbishops and bishops, prelates, doctors, clerks, and laymen obey him: also men of greatest power, riches and wisdom, together with the poor people. There is no such kingdom, and so diversely composed in the world, no not among the Mahometists. And all these verily willingly are subject to the seat: yea they have persuaded themselves that they cannot well live, that they cannot be saved, unless they be subject to the See of Rome.

And just as princes distinguish their subjects and servants by colors and signs, and common people also mark their livestock with different brands and marks, so that they may be known to belong to whom or serve whom, for every man has his colors: one white and black, another red and blue, another white and red, some black and yellow, which they give to their soldiers and servants to wear, and by this they declare themselves to be retained by him or him. And as they mark their horses with their brand and set their marks upon household vessels, so Antichrist will undoubtedly have his sign, namely, his mark, by which he may both bind people to him, and so bind them that they may be distinguished from others, and by this means wear the badge, as it were, the colors of their lord and master. And he will give his mark on the right hand, or on their foreheads.

\index{Arethas of Caesarea}\index{Primasius}According to Arethas and Primasius, and indeed all commentators agree, it explains this as the confession of the mouth, and the study and work of a good deed. We have heard how Christ, in the seventh chapter, imprinted on the foreheads of his servants the faith that works through effective charity. Indeed, the sign of God’s children is faith, and the love that arises from that faith. So do the writings of the Evangelists and Apostles testify. Nevertheless, Christ also has the external mark of his servants, those wholesome sacraments of the church, Baptism and the Lord’s Supper.

\index{Papacy}How be it if any are baptized at this day, and partake of the Lord’s Supper, call upon God the Father with the Lord’s Prayer, and utter their faith by a sincere confession of the Apostles’ Creed; moreover, confess those to be good works which are done in faith after the rule of the ten commandments, and besides this, shine in good works: shall he be taken for a true Catholic and a right Christian man? In old times, doubtless, all men would have embraced him as a brother. But what should he be at this day in the Pope’s kingdom? You shall seem, by all these things, to have confessed nothing at all of the true faith, except you plainly profess that you believe after the faith and tradition of the church of Rome; and that you do acknowledge those for good works which the church of Rome has approved. Unless you believe and profess on this wise, in vain shall you confess all the former matters. No, though you say moreover that you believe the law, the prophets, the gospel, and the Apostles. They will like you a great deal better if you say you are an obedient child of the Apostolic See and church of Rome, than if you should say that you are the child of God, a Christian man, that you put your whole trust in the Son of God, which is the only salvation and righteousness. Yea, those shall find it, who will by and by at these words cry out that they smell of heresy and a mind infected with poison. I say nothing, experience itself will witness, that I say truth. And thus does the Pope mark his men both in the forehead and right hand. Thus are the Roman whelps discerned from other faithful, as it were by marks.

Besides this, there is another thing. All Papists plainly testify that unless a man be marked in the forehead with chrism by the bishops’ hand, he is no Christian, however he be baptized and believe in Christ Jesus. Whereof it follows that they attribute more to their confirmation and anointing by the bishop than to Christian faith. Read the book called \textit{Summa angelica} in the title of confirmation. This therefore is a sealing of the Papist religion; the Christian marks of Christ are sufficient. The Pope also by another way imprints his mark on the right hand, by making of vows and performing of oaths, as they term it. For those who make a vow in entering into any religion, as they call it, do bind themselves to the Pope and See of Rome by a stipulation. Furthermore, Antichrist, the Pope, by oaths also to be performed by the holding up of the right hand, doth bind and bring into danger emperors, kings, archbishops, princes, bishops, doctors, universities, and all states. They promise that they will attempt nothing against the church of Rome, nor against the high bishop thereof, nor yet against the privileges and statutes of the See. There remain the manner of oaths in the decrees and decretals. I touch these things briefly. They see more who do not shut their eyes. And all men behold how the Pope has set his mark on the right hand and forehead of men.

\index{Turks}There follows again the cruelty and bloody tyranny which Antichrist practices against the Christians, that is to say, against those who will not receive the mark of the beast: that is, who will not prostitute themselves to the lust of the Pope and the seat of pestilence. Antichrist, says the Lord, by his power shall bring to pass that none may buy or sell, save he that hath the mark of the beast, and so forth. And these come all to one effect: the mark of the beast, the name of the beast, and the number of the name of the beast. For he hath the mark of the beast, which acknowledges the seat and professes the faith of Rome: and even to whom the Christian faith is not enough. He hath the name of the beast, whoever he be that will be named an obedient child of the holy See of Rome, and acknowledges the Pope to be head of the universal church. He hath the number of the name of the beast, which has a society with the beast, and that society that number betrays or shows. Therefore, except thou acknowledge the Pope to be supreme head of the church on Earth, with the fullness of power; unless thou profess to follow the faith of the holy church of Rome, and to detest all things whatever that See has condemned, thou art forbidden fire and water. That same has the Lord called to prohibit, that thou may neither buy nor sell. We say in Dutch, signifying one that is banished out of all men's company. He understands therefore excommunication, that horrible thunderbolt of the Pope, wherewith are stricken all those that have set more by Christ than by the Pope, or which have loathed the Pope’s decrees in comparison of the Gospel. Let him read, who desires, the Sixth Decree of Boniface the Eighth, in the Fifth Book, the Second Title, \textit{de hereticis}. Also, Clement the Fifth, Book Three, Title \textit{de hereticis}. But he that will know exactly a compendious treatise of tyranny, and a glass of butchery, let him read the Bull of Martin the Fifth, which is subject to the Sessions of the Council of Constance, and is written to Bishops and inquisitors of heretical depravity. Among other things, there is one which gives a wonderful light to this place which we now expound: where it is commanded that they do not permit those that despise the communion of the church of Rome, to keep or dwell in any house or lodging, to make any bargains, or occupy any traffic or trade of merchandise, or to have any comfort of humanity with the faithful of Christ. Read thou the rest, leaf 134. Hereunto may be added, that in Popish churches is the greatest buying and selling of all. But unless his crown be shaven, and his hands imbrewed with oil, that is, except he hath received in the forehead or head and in the right hand the mark indelible (for so they term it, that cannot be put out), he hath no merchandise left him in the house, nor so much as a little corner. But Christ whipped these merchants, or buyers and sellers, once or twice out of the temple: Antichrist hath brought them in again. And this is verily a wonder: they show more favor at this day to Jews, Turks, and heathens, than to Christians. For unto the only gospelers is no place permitted: verily for that they ascribe all to Christ, preach Christ only, leave nothing to the Pope, but rather accuse him most constantly, and bitterly.

But what shall we say to those whose hands and foreheads have been defiled with the mark? I urge them to wash themselves in the blood of Christ, forsake Antichrist, and turn to Christ, abandoning their errors and repenting of them. If you have bound yourself to Antichrist through another, do not fulfill that rash and wicked oath by ungodly speaking against the gospel. Do penance, cleanse yourself, return to Christ, and you shall be saved.

\index{Jerome, Saint}Now lest anyone here should prattle, so that we may be Christians and abundantly instructed in heavenly wisdom, even though we hear or speak nothing of the pope and papal matters, and consider those disputations unprofitable, yea odious, and pertaining to the stirring up of troubles, therefore hurtful and foolish: the Lord anticipates this, and says expressly, “Here is wisdom; in the knowledge and right judgment of these things consists the true, heavenly, and godly wisdom.” Unless we are wise in this thing, we shall be fools, and not wise. The Lord therefore excites the hearers to the study of inquiring after Antichrist, and to beware of him when he is found. For in the 14th chapter, we shall hear that they shall drink of the Wine of God’s wrath, as many as have received the mark of the beast and worshipped his image. Wherefore they shall drink at the same table with Christ of the cup of life and of the grace of God, so many as have despised papal authority. And who shall deny it to be the true wisdom, by which we may come from the wrath of God to the grace fellowship and participation of the same? Moreover, the Lord adds that those endowed with understanding, not witless and full of hurtful folly, should reckon the name of the beast; that is to say, should be diligently occupied in this matter, that those things should be diligently searched for, which worldly men affirm to be curiously sought and inquired after, [?not only without any profit, but with loss also]. Moreover, the Lord commands us to account the number of the name of the beast. He adds that the same is not hard to do, for this number to be the number of a man, to wit which a diligent man may easily by faith and industry attain to. For so does Jerome expound it also, saying that number is common and known to men. Let them leave then to trouble our godly studies, which blame our sermons made against the Pope, and laugh at our diligence such as it is in expounding papal abominations, finally which suppose we spend our time in vain in the account of times. They hear here, except they will hear nothing, that we have received commandment of the Lord so to do; moreover that the Lord testifies that wisdom is herein.

And here I give warning that the manner of speaking is to be observed, lest we and our listeners weary in vain through the insistent examination of a certain name in the numbers. For it is said to be the number of his name, as though a certain name should be gathered and composed of these characters: as most often it is gathered from these three letters, Christ. Nor do they err who think that by these three characters nothing else is signified than the name of Christ, which the Lord himself in Matthew 24 prophesied that Antichrist should use. Surely, he calls himself Christ’s vicar. I know very well that the proper names of great men have been sealed by prophecies and signified before: as Josiah, Cyrus, Jesus. But here you can gather no such thing, but forcibly and as it were against the text. I understand therefore by the number of the name of Antichrist or beast, the very acceptance whereby we come unto his name. And a name is a brief definition or description of any thing, whereby it is known of what sort and manner it is. Which thing in this our cause, the nobility supplies, which brings us unto those times which give him his name, from which he takes his name, that is which times reveal unto us Antichrist spoken of before in the prophets, and show us who and what he is, or who we should take for Antichrist, even him verily, which having brought low three kings, he himself starts up of naught, and to the destruction of the true religion begins to reign.

\index{Apocalypse (Book of)}\index{Irenaeus}And now he shows us expressly this number, which I may call nominal, and vocal, which may lead us to Antichrist, that we may know who it is; and when we know him, we may beware of him, and commands us to number the years six hundred sixty and six. For so many signify these Greek letters. In explaining this number, commentators have varied greatly. I prefer the explanation of the blessed Martyr Irenaeus, who, perhaps a hundred years after the publication of the Apocalypse, wrote his book against heresies, and encountered some who had heard John preach. Irenaeus also includes Andreas, the good bishop of Caesarea, who, with Arethas, speaks thus: the perfect reckoning and just account of the noble [?], as likewise other things which are written of the same Antichrist, the opportunity of time will reveal, and experience will confirm, to those who watch diligently. For if it were necessary, as some of the Doctors suppose, that this name should be manifestly known, he who saw him would surely have revealed it. But divine grace did not allow the name of this pestiferous beast to be rehearsed in this godly book. Thus far Andreas.

After the same manner, the holy Martyr of Christ Irenaeus, before Andreas, wrote in the fifth book against Heresies. For at the end of the book, he says it is safer and without danger to wait for the fulfillment of the prophecy than to suspect and guess at every name, since many names may be found having the aforementioned number, whereby the question is not answered. Yet, he adds immediately, “the name contains the number 666.” And it is very likely to be true. For this term has a kingdom. For the Latins now reign. This is what he says.

And without doubt, this good doctor did not err in the slightest, being endowed with the Holy Spirit of God. For we see that the church of Rome is called the Latin church, and the Pope the high bishop of the Latin church. We see all services in churches performed in the Latin tongue. In courts and all judgments of bishops, the Latin tongue is used alone. Moreover, no one may serve in this church unless he is a Latinist. What will you say that these Latinists call the Hebrew, that is to say, the holy tongue, by a contemptuous name, “Jewish”? They call the Greek church and tongue heretical. The Bibles in Greek and Hebrew are suspected by them. For they insist that only the Latin Bibles are authentic and should be read by all as authentic. But these things are better known than that I need to admonish and recite them here with many words.

Nevertheless, this holy man Irenaeus does not wholly affirm this conjecture as most certain, although what he said was most probable and likely to be true. For he adds, “we will not be in danger herein” (for he also recited the name [illegible], the royal or tyrannical name of Nimrod), “neither will we definitively pronounce that he shall have this name,” knowing that if it were necessary for his name to be manifestly preached at this present time, it would certainly have been uttered by him who saw the Revelation. But he showed the number of the name that we might beware of him when he comes, knowing who he is. And he concealed his name, for it is not worthy to be preached by the Holy Spirit and so forth.

Nevertheless, this same matter shows us how to understand those 666 years. For he says: knowing the certain number, which Scripture reveals to us—that is, the number 666—let the faithful remain steadfast or look for this: first, the division of the kingdom into ten, then the same power reigning and beginning to reform its affairs and to expand its kingdom. After this, one who comes suddenly will challenge a kingdom for himself and will put the aforementioned kings into fear, having a name containing that same number, so that he may be recognized certainly as the abomination of desolation. He says this again.

\index{Daniel (Book of)}But who does not see that the holy Martyr [Jerome] sends us to the prophecy of Daniel, which in chapter 7 says how the Roman Empire shall be divided into many kingdoms; and how in the midst of those kings a little horn should arise, which would overthrow and abase three horns; and that the same would begin to reign proudly, tyrannically, and wickedly, against both God and men, but chiefly to the faithful, intolerably?

Let us then see how and when these things are fulfilled. Where the Roman Empire had godly emperors, even then wicked Rome would not bow her stiff neck to Christ, but always most obstinately aspired to her old and wonted idolatry, which it coveted to have restored. And finally, when the fatal time was at hand, wherein the Lord, most righteous, thought to requite bloody Rome, he armed against her the Goths, Vandals, and Germans, which subdued and destroyed the lady of the whole world, and destroyed the whole empire. On this matter, seek more in sermon 57, and in the sermons following.

And it is evident by Histories that the Roman Empire, the Goths beginning to invade it, did decline, with provinces revolting in every place, and was divided into many kingdoms. For to speak nothing of Asia and Africa, the Persians ravaged the former, and the Vandals the latter, all Greece followed the Emperor of Constantinople, and likewise other nations never. The Westgoths possessed all of Spain, and the Franks subdued Gaul, Germany, and the nations adjoining to the same. The Eastgoths and Lombards obtained Italy. Thus, truly, many kingdoms were established, and in place of Rome reigned many kings. However, while these kings considered how they might best enlarge their kingdoms and put down and expel others, full craftily the Bishop of Rome played his part also. For he obtained supremacy over all bishops and so gained great authority with kings and realms, yea, and linked himself in league and amity with kings and princes. Whereupon quickly and suddenly, or as the Martyr of Christ prophesied, upon the sudden, he appeared, and at last usurped a kingdom, to wit, of Rome. For by his judgments falsely taken for apostolic, he put down King Childeric, of the lineage of Merovinges, the lawful king of France, and advanced Pipin, then Captain of the French guard, to the crown. And so he overthrew or plucked down one horn and bound unto him a most mighty king, by whose power afterward he was a terror to the kings of Greece and Lombardy.

Concerning the year of our Lord 269, the Emperor of Constantinople, expelling the East Goths, instituted a new government in Italy. But since this kind of rule and governance is not known to all men, I will briefly recite what and how great it was, by the words of Nauclerus the Historian, from his \textit{Generations}, book 20. “Then began the City of Rome and Italy to have a new manner of governance, by which they lost more of the dignity, glory, and fear over all the world than of all the calamities which these 160 years have afflicted them, and at the last had left Rome to be inhabited of wild beasts.” For Longinus brought in a new name of dignity, the \textit{exarchate} of Italy: that is, the high magistrate. Which keeping still at Ravenna, went never to the City of Rome. And in the government of Italy and of Cities, he kept first this order, that the president should not govern the province or region, but every City had their magistrate to govern them, whom he called Dukes. Wherefore making Rome equal with other Cities and Towns, in this thing only he honored the same, that he called the magistrate’s place in Rome, \textit{praesides}. But they that did succeed him were called Dukes, as they were afterward for many years, so that it was called the Dukedom of Rome, as the Dukedom of Narnia and Spolet. Neither after Narses and Basil had it any more either Consuls, or Senate lawfully assembled: but by a Duke of Greece, whom the high magistrate sent from Ravenna, the commonwealth of Rome was governed a long time.” Thus much he.

Any man may easily perceive here that the prophecies are fully accomplished, and the Roman Empire has fallen into ashes. For she who had been the most powerful ruler in the world is now seen to be a vile government, no more excellent than that of Spoleto and Narnia. And it is to be known that this Exarchate in Italy was the third lordship instituted since Augustulus was slain, in whom the histories say that the Empire of the west was finished and ended. For first, when Augustulus was slain, the Germans under their king Odacer possessed Rome. Afterward, the Ostrogoths, under the leadership of their Duke Theodoric of Verona, expelled and killed Odacer and reigned at Rome and in Italy. Last of all, the Ostrogoths were expelled and killed by the Lombards, and this Exarchate was instituted. And the Lombards, being called into Italy by the Greeks against the Goths, would no longer leave, for they saw the land fertile and rich, pleasant and abundant with pleasures. Possessing great power in Italy, they subdued many cities and people of Italy, establishing now the fourth dominion, which they called the kingdom of the Lombards. They had powerful kings. However, that Exarchate of Ravenna, although they waited diligently for it and attempted to invade it, could never extinguish it, until the Bishop of Rome put his hand to it, pretending the sincerity of religion.

Historians record sixteen Exarchs in succession, who ruled for approximately one hundred and eighty years. The fifteenth was named Paulus. Nauclerus states that in January of the twenty-fifth year, Leo the Third, Emperor of Constantinople, ordered that all those subject to the Roman Empire should tear down their images, break them, and burn them. Conversely, the Pope—some say Gregory the Second, others Gregory the Third—wrote to the entire world, instructing them not to obey these wicked commands of the Emperor. Platina writes more in the life of Gregory the Third: Gregory, with the consent of the clergy of Rome, deposed Leo the Third, Emperor of Constantinople, from both the Empire and the communion of the faithful, because he had torn down images. Nauclerus adds that the papal decrees held such great authority at that time that first the people of Ravenna, and then the people and soldiers of Venice, openly rebelled against the Emperor and the Exarch in Italy. And the treason increased daily. For Marinus Spatarius, Duke of Rome, and his son Adriane, while traveling through Campania, were slain by the Romans. In their place, they created one Peter as Duke of Rome. The people of Ravenna also, while some sided with the Emperor and others with the Pope, in a riot, killed Paulus the Exarch and his son. Thus writes Nauclerus.

In these disturbances, the Lombards, believing the opportunity they had long desired had arrived, invade the lands of the empire through the agency of their king, Luitprand, and besiege Rome itself. But Pope Gregory, the instigator of all the unrest in Italy, the soldier and practitioner of it, and resembling neither a priest nor a preacher, summons Charles Martel, father of King Pippin, with his French warriors into Italy against the Lombards. However, this Charles persuades the king of the Lombards amicably to withdraw from the city. But not long afterward, Aistulf, king of the Lombards, plunders again the lands of Ravenna, renews the Italian war, and seizes Ravenna itself, demanding tribute from the city of Rome. But Pope Stephen II, who aspired to the governance of Ravenna and wished for the Lombards to be destroyed, requests aid from King Pippin of France, to whom, not long since, Pope Zacharias by his unjust judgment [? a false decree] (as many suppose) had given the kingdom, essentially offering him the kingdom. Therefore, the Frenchmen, armored and eager to conquer Italy, are present. While King Pippin entered Italy, he encountered the ambassadors of the Emperor of Constantinople, who requested that he restore Ravenna, the exarchate, and its lands to the empire, which rightfully belonged to it, and not to the Pope or the Romans. Pippin replied that he was warring for Saint Peter and the Pope, and to prevent the Lombards from harassing the church, and that he would take the exarchate and other territories of Italy from them and deliver them to the Pope, which he did. For he overcame King Aistulf, took from him the governance of Ravenna, and delivered it to the Bishop of Rome.

Herein may all men see, unless it be those who will see nothing, how this contemptuous bishop, and very small horn, has at one push overthrown two horns. For he has put the emperor of Constantinople from the government of Italy and has put down the King of Lombardy, causing his people to be driven out of Italy. For a few years after, the Pope, by the force of Charlemagne, put down Desiderius, the last king of Lombardy, and destroyed withal the whole people of the Lombards. And thus the Pope arose and became as it were king of old Rome and of the chief part of Italy. And now are the beginnings of the kingdom laid, but as yet he reigned not with full authority, as is declared before. Eberardus, Bishop of Salisbury, whose words I recited in the preface of this book, extends these things further. But I suppose this our exposition to accord with the prophet, with the things and times. And the Pope gave to King Pippin for so great a donation a title, as Platina shows in the life of Stephen the second, that all kings of France should be called most Christian. Afterward was the Image of the Empire bestowed upon Charles, whereof is spoken before.

And lest the Pope should seem to have received nothing, while King Pippin gave him the exarchate, the stories report thus: the exarchate was divided into two regions, in Pentapolis and Aemilia. Pentapolis had five cities: Ravenna, Cesena, Classe, Forum Livii, and Forum Popilii. In Aemilia were Bologna, Reggio, Parma, Piacenza, and all the lands that lie from the borders of the Piacentines and Ticini to Adria, and from Adria to Ariminum, \&c. But he who desires may read the Donation of Louis the Pious, in Volaterranus’s Geography, where he names the kings of France. We say nothing yet of how he afterward usurped power to himself over kings and realms, finally over all churches and souls, so that we must confess that a more marvelous prince never lived.

You have here a brief and concise account, declaring how the Pope, having humbled and overthrown three kings, began to be made a king himself. But let us now apply the number of the name of the beast, so that it may be known to the whole world, that there is no other Antichrist to be looked for than the bishop of Rome, who has indeed laid the foundation of his kingdom under the emperor Phocas, built it under the kings of France, and enlarged it under the emperors Henry and Frederick, finally having established it under the emperors following; reigns in our time, and has done so for certain ages already past, \&c.

\index{Eusebius of Caesarea}The calculation of 666 years must be reckoned from the time when John saw the Apocalypse. Irenaeus says it was seen not long since, but in a manner in our days, about the end of the reign of Domitian. And Eusebius in his chronicles says that it was in the year of our Lord 97. Therefore, there remain yet three years to complete a hundred years from the birth of our Lord. Add to a hundred years these years of the number of the beast, 666, and subtract those three years of the first hundred, and you shall have the year of our Lord 763, which was the 13th year, or thereabouts, of King Pepin’s reign, and the 7th of Pope Paul. Notwithstanding that there be writers of histories and times who attribute to Paul but one year, etc.

Now must we not look only at what things happened in the very year 763, but what followed in the years before and after. Of this, I will recite a few things from the writers of histories and chronicles.

Nauclerus, in the sixteenth generation, says that in the year of our Lord 1150, under Pope Zacharias and Emperor Constantine V, the twenty-sixth generation began, in which there was an alteration of the Kingdom of France, an abolishment of the Kings of Lombardy, and a translation of the Roman Empire from the Greeks. These great alterations were happily foreshadowed by the wonders that occurred at this time. In Mesopotamia, the earth split open for two miles, and a mule was said to have spoken with a man’s voice. Ashes fell down from heaven. There were wonderful earthquakes. Crosses appeared upon men’s garments. These things Nauclerus wrote. Similar events are recorded in the history of Eutropius, in the twenty-second book, under the year of Constantine VI; moreover, in Vincent’s historical chronicle, and in the Fasciculus Temporum.

In the year of our Lord 751, through the counsel of Zachary the Pope, Pipin, master of the king’s household, oppressed his lord Hilderych, king of France, began to reign, and reigns. Eighteen years. This writes Aemilius in the second book of kings of France. And in the year 755, Pipin enters Italy with an army, vanquishes the king of Lombards, and gives the entire government of Ravenna to Saint Peter, against the will of the Emperor of Constantinople. Vespergenensis in chronicles. You see how, instead of the Emperor, the Pope begins, in a way, to reign at Rome and in Italy, the horns are shaken off, according to the prophecy. Matthew Palmer in his chronicles, under the year 756, says that the Roman Empire revolted at a pace in the East, and the Emperor was persecuting the Christians (Idolaters he should have termed them). Pope Stephen gave to the kings of France the imperial titles and dignities, and confirmed Pipin and the successors of his lineage as kings only, all others utterly excluded, and in the name of the people of Rome, called him Patrician. Hitherto Palmer.

John Functius in his Chronicle writes: In the year of our Lord 756, the rites and ceremonies of the church of Rome were carried into France and first received. In the year of our Lord 757, Paul is made Pope, and immediately follows that fatal year of our Lord 763, which falls as the midpoint between the year 750 and 770, or 773. Within this time, all those things occurred that both give the name to Antichrist, and whereby, just as everything else is known by his name, so too he receives his name and is known.

In the year of our Lord 768, Stephen III held a council at Rome in the Lateran church with the bishops of France and Italy: and decreed that no one should be ordained bishop of Rome except a Cardinal. He also condemned the Greek council of the Emperor Constantine against images, which he commanded to be possessed and worshipped. These things writes Antoninus in Chronicles, title 14, chapter 1 and 5.

After this, that great Charles, son of Pipin, was summoned to Italy by Pope Adrian, took Desiderius, King of the Lombards, and extinguished the Kingdom of the Lombards. This occurred in the year of our Lord 773, and also the two hundred and fourth year after the Lombards arrived in Italy. He confirmed and augmented the donation of Pipin his father, as many historians relate. John Functius in his Chronicle adds that throughout the entire realm of France, at the command of Charles, the ceremonies of the Roman church were instituted. Now we have the name of Antichrist, of the number 666. We know who he is, and whom we should beware of.

I cannot here omit, but must note a few words concerning the Sibyl’s prediction regarding the origin of the Antichrist, which, in my judgment, is very consistent with what has been said before. For the eight books of the Sibyl’s oracles, taken from the library of the honorable commonwealth of Augsburg, were published by the most godly and learned Doctor Sixtus Betuleius in the year of our Lord 1545, and that in Greek. And this Sibyl of Erythrea, or whatever she was, prophesies in the eighth book that Rome shall fall and be burned with fire. The words of the Sibyl in Greek are to this effect.

And this event occurred at the time Totila, King of the Goths, burned the city, as we have recounted previously. Shortly thereafter, these words are added in the same prophetic writings: When emperors have oppressed the world with great bondage from east to west, the number fifteen will be fulfilled. Then a king will appear wearing a white hat, who will never add his name to Pontius [?], as if he were to be called Pontifex. He will live for worldly pleasures and will reward with his wicked foot. And the rest that are read there.

She bids us count fifteen kings from the burning of Rome. After whom a new king will come, whom she describes. And it is manifest that Rome was taken, plundered, and burned under Emperor Justinian. After Justinian the Younger, fifteen kings are counted to Emperor Theodosius. After Theodosius, Leo III succeeds, whom she calls “delicatos,” that is, given to pleasures. Because most of them were not very valiant, but under Leo III, Italy revolts from the emperor. And shortly also the government called the Exarchate was given to the pope by King Pippin, against the emperor’s mind. We see therefore that the reckonings agree. For we have also brought the years to King Pippin, 666. And so a new king arises, whom the Sibyl names [illegible] notable because of his white hat or mitre. For so she notes the bishop (who in old time wore white miters on their heads) that should be a king. She gives him a name also. For she says how he has a name never unto Ponti. For add to the word Ponti, \textit{fex}, and you have Pontifex. She annexes certain notes or marks also: that he shall regard earthly things, and not heavenly; and that he shall also provide [illegible] and give rewards with his ungracious foot. And that is rightly spoken, since after Domitian and Diocletian, none of all the kings, save the Pope, has offered his foot to be kissed: whereby fools think they receive great rewards. But omitting these things, let us return into the way.

The blessed martyr Irenaeus, speaking of this king, in the same fifth book, says that in the beast’s coming, there is a recapitulation of all iniquity and all deceit, so that all apostatical power concurring and concluded in him might be thrown into a furnace of fire. And that he has spoken this thing by the spirit of prophecy, all men will confess who have read the lives of the Bishops of Rome: but especially of Sylvester II, Benedict IX, Gregory VI, Gregory VII, Urbanus II, Paschalis II, Alexander III, Innocentius III, Gregory IX, Boniface VIII, Clement V, John XXII, to speak nothing of various others. What has been done in our days by Julius, Clement, Leo, and Paul? Spain, France, England, Hungary, and Germany, and other realms speak, which have been set together by the ears and entangled among themselves with most cruel wars. The blood of martyrs shed speaks, which cries unto the Lord. What remains therefore, but that we should take heed to ourselves and beware of this man of sin, and cleave to our Redeemer Christ our Lord, beseeching him that he would come shortly and deliver us from all evil. Amen, Amen.
